\begin{center}
\textbf{{\Large ABSTRACT}}\\
\begin{center}
\justifying
Linear recurrent networks attain inference cost linear in sequence length, but have historically been weak at state tracking, the task of maintaining a structured hidden state whose correct value depends on the entire preceding input. Recent work closes much of this gap by constructing each state-transition matrix as a product of $n_h$ generalized Householder factors, where a larger $n_h$ yields a richer set of reachable transitions at proportionally greater computational cost. In existing formulations $n_h$ is a single global hyperparameter, fixed before training and applied identically to every token.

This work establishes that the number of Householder factors a symbol requires is not a property of that symbol. It is determined by the algebraic closure of the entire input alphabet, through the minimum dimension of a faithful representation of the group that the alphabet generates. The resulting order-requirement model is verified by explicit construction in double precision: a five-cycle requires four factors when the alphabet generates the symmetric group $S_5$, and only two when the alphabet generates the alternating group $A_5$, which admits a three-dimensional faithful representation. A direct consequence is that introducing a single transposition, the least expressive symbol available, doubles the requirement of every five-cycle in the alphabet. Two hard bounds on reachability are also established: a rank bound that forces the factor count to be at least the transposition length of the target permutation, and a determinant-parity bound that governs the admissible padding of shorter solutions.

Building on this analysis, an adaptive-order layer is proposed in which a monotone halting gate assigns a per-token factor count, trained end to end with a penalty on expected computation and no supervision on the allocation itself. Under an oracle allocation, in which the correct per-token order is supplied from ground truth, the layer attains accuracy above $0.9997$ at $1.15$ to $1.75$ factors per token on alphabets where the cheapest adequate fixed order costs four, a reduction in required work of $2.29\times$ to $3.48\times$. Trained end to end across three random seeds, the halting gate reproduced the oracle allocation exactly in one seed, assigning one factor to every low-requirement token and four to every high-requirement token over $256{,}000$ evaluated tokens with zero mean absolute allocation error. A supporting stability analysis over fifty training runs establishes that outcomes on the harder alphabets are bimodal rather than concentrated, and that the configurations in which adaptive order offers the largest reduction are precisely those that train reliably.
\end{center}
\end{center}
\keywords{Linear recurrent networks, state tracking, Householder transformations, adaptive computation, group representations, conditional computation}\\
